% Part: first-order-logic
% Chapter: completeness
% Section: construction-of-model

\documentclass[../../../include/open-logic-section]{subfiles}

\begin{document}

\iftag{FOL}
      {\olfileid[pt]{fol}{com}{mod}}
      {\olfileid[pt]{pl}{com}{mod}}

\olsection{Construção de um Modelo}

\begin{explain}
\iftag{FOL}{Por ora, não nos preocupamos com $\eq$, isto é, queremos apenas
  mostrar que um conjunto consistente~$\Gamma$ de !!{sentence}s que não
  contém~$\eq$ é satisfatível. Primeiro estendemos~$\Gamma$ a um conjunto
  consistente, !!{complete} e saturado~$\Gamma^*$.  Nesse
  caso, a definição de um modelo~$\Struct{M(\Gamma^*)}$ é simples: Tomamos
  o conjunto de termos fechados de~$\Lang{L'}$ como o domínio. Atribuímos
  cada !!{constant} a si mesmo e nos asseguramos de que, mais geralmente, para
  todo termo fechado~$t$, $\Value{t}{M(\Gamma^*)} = t$.  Aos
  !!{predicate}s são atribuídas extensões de tal forma que uma !!{sentence}
  atômica é verdadeira em $\Struct{M(\Gamma^*)}$ se e somente se está
  em~$\Gamma^*$.  Isso obviamente tornará verdadeiras em $\Struct{M(\Gamma^*)}$ todas as
  !!{sentence}s atômicas em~$\Gamma^*$. As demais
  são verdadeiras desde que o $\Gamma^*$ com que começamos seja consistente,
  completo e saturado.}{Estamos agora prontos para definir !!a{valuation}
  que torna verdadeiras todas as $!A \in \Gamma$. Para fazer isso, primeiro aplicamos o
  Lema de Lindenbaum: obtemos um $\Gamma^* \supseteq
  \Gamma$ consistente completo. Deixamos as !!{propositional variable}s em~$\Gamma^*$
  determinarem~$\pAssign v(\Gamma^*)$.}
\end{explain}

\iftag{FOL}{%
\begin{defn}[Modelo de termos]
\ollabel{defn:termmodel}
Seja $\Gamma^*$ um conjunto consistente, !!{complete} e
saturado de !!{sentence}s em uma linguagem~$\Lang L$. O \emph{modelo de
  termos}~$\Struct M(\Gamma^*)$ de $\Gamma^*$ é a !!{structure}
definida como segue:
\begin{enumerate}
\item O !!{domain}~$\Domain{M(\Gamma^*)}$ é o conjunto de todos os termos
  fechados de~$\Lang L$.
\item A interpretação de !!a{constant} $c$ é o próprio $c$:
  $\Assign{c}{M(\Gamma^*)} = c$.
\item Ao !!{function}~$f$ é atribuída a função que, dados como
  argumentos os termos fechados $t_1$, \dots, $t_n$, tem como valor o
  termo fechado $f(t_1, \dots, t_n)$:
\[
\Assign{f}{M(\Gamma^*)}(t_1, \dots, t_n) = f(t_1,\dots, t_n)
\]
\item Se $R$ é um !!{predicate} de $n$ lugares, então
  \[
  \tuple{t_1, \dots,
  t_n} \in \Assign{R}{M(\Gamma^*)} \text{ se e somente se } \Atom{R}{t_1, \dots,
    t_n} \in \Gamma^*.
  \]
\end{enumerate}
\end{defn}
}{%
\begin{defn}
\ollabel{defn:termmodel}
Suponha que $\Gamma^*$ seja um conjunto consistente completo de !!{formula}s. Então deixamos
\[
\pAssign v(\Gamma^*)(p) = \begin{cases}
  \True & \text{se $p \in \Gamma^*$}\\
  \False & \text{se $p \notin \Gamma^*$}
  \end{cases}
\]
\end{defn}
}

\iftag{FOL}{%
Verificaremos agora que de fato temos $\Value{t}{M(\Gamma^*)} = t$.

\begin{lem} 
\ollabel{lem:val-in-termmodel} Seja $\Struct M(\Gamma^*)$ o modelo de termos
de \olref{defn:termmodel}; então $\Value{t}{M(\Gamma^*)} = t$.
\end{lem}
\begin{proof} 
 A prova é por indução sobre $t$, onde o caso base, quando $t$
 é !!a{constant}, segue diretamente da definição do modelo de
 termos. Para o passo de indução, assuma que $t_1, \ldots, t_n$ são termos fechados
 tais que $\Value{t_i}{M(\Gamma^*)} = t_i$ e que $f$ é !!a{function}
 $n$-ária. Então 
\begin{align*}
\Value{f(t_1,\ldots,t_n)}{M(\Gamma^*)} &= \Assign{f}{M(\Gamma^*)}(\Value{t_1}
 {M(\Gamma^*)}, 
\ldots, \Value{t_n}{M(\Gamma^*)}) \\
&= \Assign{f}{M(\Gamma^*)}(t_1, \dots, t_n) \\
&= f(t_1,\dots, t_n),
\end{align*} 
e, assim, por indução, isto vale para todo termo fechado~$t$.
\end{proof}
}

\iftag{FOL}{%
\begin{explain}
  \iftag{prvEx}{!!^a{structure}~$\Struct{M}$ pode tornar verdadeira uma
    !!{sentence}~$\lexists[x][!A(x)]$ existencialmente quantificada sem que
    haja uma instância~$!A(t)$ que ela torne verdadeira.}{}
  \iftag{prvAll}{!!^a{structure}~$\Struct{M}$ pode tornar verdadeiras todas
    as instâncias~$!A(t)$ de uma !!{sentence}~$\lforall[x][!A(x)]$
    universalmente quantificada sem
    tornar~$\lforall[x][!A(x)]$ verdadeira.}{} Isso ocorre porque, em geral,
  nem todo !!{element} de~$\Domain{M}$ é o valor de um termo fechado
  ($\Struct{M}$ pode não ser coberta).  Esta é a razão pela qual a
  relação de satisfação é definida por meio de atribuições de variáveis. Entretanto,
  para o nosso modelo de termos~$\Struct{M(\Gamma^*)}$ isso não seria
  necessário --- porque ela é coberta. Este é o conteúdo do próximo
  resultado.
\end{explain}

\begin{prop}
\ollabel{prop:quant-termmodel}
Seja $\Struct M(\Gamma^*)$ o modelo de termos de \olref{defn:termmodel}.
\begin{tagenumerate}{prvEx,prvAll}
\tagitem{prvEx}{$\Sat{M(\Gamma^*)}{\lexists[x][!A(x)]}$ se e somente se
  $\Sat{M(\Gamma^*)}{!A(t)}$ para pelo menos um termo fechado~$t$.}{}
\tagitem{prvAll}{$\Sat{M(\Gamma^*)}{\lforall[x][!A(x)]}$ se e somente se
  $\Sat{M(\Gamma^*)}{!A(t)}$ para todo termo fechado~$t$.}{}
\end{tagenumerate}
\end{prop}

\begin{proof}
\begin{tagenumerate}{prvEx,prvAll}
\tagitem{prvEx}{%
  \iftag{probEx}{Exercício.}{Por \olref[syn][ass]{prop:sat-quant},
    $\Sat{M(\Gamma^*)}{\lexists[x][!A(x)]}$ se e somente se, para pelo menos uma
    atribuição de variáveis~$s$, $\Sat{M(\Gamma^*)}{!A(x)}[s]$. Como
    $\Domain{M(\Gamma^*)}$ consiste nos termos fechados de~$\Lang{L}$,
    este é o caso se e somente se há pelo menos um termo fechado~$t$ tal
    que $s(x) = t$ e $\Sat{M(\Gamma^*)}{!A(x)}[s]$.  Por
    \olref[fol][syn][ext]{prop:ext-formulas},
    $\Sat{M(\Gamma^*)}{!A(x)}[s]$ se e somente se $\Sat{M(\Gamma^*)}{!A(t)}[s]$,
    onde $s(x) = t$.  Por \olref[fol][syn][ass]{prop:sentence-sat-true},
    $\Sat{M(\Gamma^*)}{!A(t)}[s]$ se e somente se $\Sat{M(\Gamma^*)}{!A(t)}$,
    já que $!A(t)$ é uma sentença.}}{}
\tagitem{prvAll}{%
  \iftag{probAll}{Exercício.}{Por \olref[syn][ass]{prop:sat-quant},
    $\Sat{M(\Gamma^*)}{\lforall[x][!A(x)]}$ se e somente se, para toda
    atribuição de variáveis $s$, $\Sat{M(\Gamma^*)}{!A(x)}[s]$. Lembre que
    $\Domain{M(\Gamma^*)}$ consiste nos termos fechados de~$\Lang{L}$,
    de modo que, para todo termo fechado~$t$, $s(x) = t$ é uma tal atribuição de
    variáveis, e, para qualquer atribuição de variáveis, $s(x)$ é algum termo
    fechado~$t$.  Por \olref[fol][syn][ext]{prop:ext-formulas},
    $\Sat{M(\Gamma^*)}{!A(x)}[s]$ se e somente se $\Sat{M(\Gamma^*)}{!A(t)}[s]$,
    onde $s(x) = t$.  Por \olref[fol][syn][ass]{prop:sentence-sat-true},
    $\Sat{M(\Gamma^*)}{!A(t)}[s]$ se e somente se $\Sat{M(\Gamma^*)}{!A(t)}$,
    já que $!A(t)$ é uma sentença.}}{}
\end{tagenumerate}
\end{proof}
}{}

\tagprob[FOL]{probEx,probAll}
\begin{prob}
Complete a prova de \olref[fol][com][mod]{prop:quant-termmodel}.
\end{prob}
\tagendprob

\begin{lem}[Lema da Verdade]
\ollabel{lem:truth} \iftag{FOL}{Suponha que $!A$ não contenha~$\eq$. Então}{}
$\iftag{FOL}{\Sat{M(\Gamma^*)}}{\pSat{v(\Gamma^*)}}{!A}$ se e somente se $!A \in \Gamma^*$.
\end{lem}

\begin{proof}
Provamos ambas as direções simultaneamente e por indução sobre $!A$.
\begin{enumerate}
\tagitem{prvFalse}{\indcase{!A}{\lfalse}
  {$\iftag{FOL}{\Sat/{M(\Gamma^*)}{\lfalse}}{\pSat/{v(\Gamma^*)}{\lfalse}}$
    por definição de satisfação. Por outro lado, $\lfalse \notin
    \Gamma^*$ já que $\Gamma^*$ é consistente.}}{}

\tagitem{prvTrue}{\indcase{!A}{\ltrue}
  {$\iftag{FOL}{\Sat{M(\Gamma^*)}}{\pSat{v(\Gamma^*)}}{\ltrue}$
    por definição de satisfação. Por outro lado, $\ltrue \in
    \Gamma^*$ já que $\Gamma^*$ é consistente e !!{complete}, e
    $\Gamma^* \Proves \ltrue$.}}{}

\tagitem{FOL}{\indcase{!A}{R(t_1, \dots, t_n)}
  {$\Sat{M(\Gamma^*)}{\Atom{R}{t_1, \dots, t_n}}$ se e somente se $\tuple{t_1,
      \dots, t_n} \in \Assign{R}{M(\Gamma^*)}$ (pela definição de
    satisfação) se e somente se $R(t_1, \dots, t_n) \in \Gamma^*$ (pela
    construção de $\Struct
    M(\Gamma^*)$).}}{\indcase{!A}{p}{$\pSat{v(\Gamma^*)}{p}$ se e somente se
    $\pAssign v(\Gamma^*)(p) = \True$ (pela definição de
    satisfação) se e somente se $p \in \Gamma^*$ (pela construção de
    $\pAssign v(\Gamma^*)$).}}

\tagitem{prvNot}{%
  \indcase{!A}{\lnot !B}
    {$\iftag{FOL}{\Sat{M(\Gamma^*)}}{\pSat{v(\Gamma^*)}}{\indfrm}$ se e somente se
    $\iftag{FOL}{\Sat/{M(\Gamma^*)}{!B}}{\pSat/{v(\Gamma^*)}{!B}}$ (por
    definição de satisfação). Por hipótese de indução,
    $\iftag{FOL}{\Sat/{M(\Gamma^*)}{!B}}{\pSat/{v(\Gamma^*)}{!B}}$ se e somente se
    $!B \notin \Gamma^*$. Como $\Gamma^*$ é consistente e
    !!{complete}, $!B
    \notin \Gamma^*$ se e somente se $\lnot !B \in \Gamma^*$.}}{}

\tagitem{prvAnd}{%
\iftag{probAnd}{%
  \indcase!{!A}{!B \land !C}{}}{%
  \indcase{!A}{!B \land
    !C}{$\iftag{FOL}{\Sat{M(\Gamma^*)}}{\pSat{v(\Gamma^*)}}{\indfrm}$
    se e somente se temos ambos
    $\iftag{FOL}{\Sat{M(\Gamma^*)}}{\pSat{v(\Gamma^*)}}{!B}$ e
    $\iftag{FOL}{\Sat{M(\Gamma^*)}}{\pSat{v(\Gamma^*)}}{!C}$ (por
    definição de satisfação) se e somente se $!B \in \Gamma^*$ e $!C \in
    \Gamma^*$ (pela hipótese de indução). Por
    \olref[ccs]{prop:ccs}\olref[ccs]{prop:ccs-and}, este é o caso
    se e somente se $(!B \land !C) \in \Gamma^*$.}}}{}

\tagitem{prvOr}{%
\iftag{probOr}{%
  \indcase!{!A}{!B \lor !C}{}}{%
  \indcase{!A}{!B \lor !C}
    {$\iftag{FOL}{\Sat{M(\Gamma^*)}}{\pSat{v(\Gamma^*)}}{\indfrm}$
    se e somente se $\iftag{FOL}{\Sat{M(\Gamma^*)}}{\pSat{v(\Gamma^*)}}{!B}$ ou
    $\iftag{FOL}{\Sat{M(\Gamma^*)}}{\pSat{v(\Gamma^*)}}{!C}$ (por
    definição de satisfação) se e somente se $!B \in \Gamma^*$ ou $!C \in
    \Gamma^*$ (por hipótese de indução). Este é o caso se e somente se $(!B
    \lor !C) \in \Gamma^*$ (por
    \olref[ccs]{prop:ccs}\olref[ccs]{prop:ccs-or}).}}}{}
  
\tagitem{prvIf}{%
\iftag{probIf}{%
  \indcase!{!A}{!B \lif !C}{}}{%
  \indcase{!A}{!B \lif !C}
    {$\iftag{FOL}{\Sat{M(\Gamma^*)}}{\pSat{v(\Gamma^*)}}{\indfrm}$
    se e somente se $\iftag{FOL}{\Sat/{M(\Gamma^*)}}{\pSat/{v(\Gamma^*)}}{!B}$ ou $\iftag{FOL}{\Sat{M (\Gamma^*)}}{\pSat{v(\Gamma^*)}}{!C}$ (por
    definição de satisfação) se e somente se $!B \notin \Gamma^*$ ou $!C \in
    \Gamma^*$ (por hipótese de indução). Este é o caso se e somente se $(!B
    \lif !C) \in \Gamma^*$ (por
    \olref[ccs]{prop:ccs}\olref[ccs]{prop:ccs-if}).}}}{}

\iftag{FOL}{%
\tagitem{prvAll}{%
  \iftag{probAll}{%
    \indcase!{!A}{\lforall[x][!B(x)]}{}}{%
    \indcase{!A}{\lforall[x][!B(x)]}{$\Sat{M(\Gamma^*)}{\indfrm}$ se e somente se
      $\Sat{M(\Gamma^*)}{!B(t)}$ para todos os termos~$t$
      (\olref{prop:quant-termmodel}).  Por hipótese de indução, este
      é o caso se e somente se $!B(t) \in \Gamma^*$ para todos os termos~$t$; por
      \olref[hen]{prop:saturated-instances}, isto, por sua vez, é o caso
      se e somente se $\lforall[x][!A(x)] \in \Gamma^*$.}}}{}

\tagitem{prvEx}{%
  \iftag{probEx}{%
    \indcase!{!A}{\lexists[x][!B(x)]}{}}{%
    \indcase{!A}{\lexists[x][!B(x)]}{$\Sat{M(\Gamma^*)}{\indfrm}$ se e somente se
      $\Sat{M(\Gamma^*)}{!B(t)}$ para pelo menos um termo~$t$
      (\olref{prop:quant-termmodel}).  Por hipótese de indução, este
      é o caso se e somente se $!B(t) \in \Gamma^*$ para pelo menos um termo~$t$.
      Por \olref[hen]{prop:saturated-instances}, isto, por sua vez, é o
      caso se e somente se $\lexists[x][!B(x)] \in \Gamma^*$.}}}{}
}{}
\end{enumerate}
\end{proof}


\tagprob[FOL]{probOr,probAnd,probIf,probEx,probAll}

\begin{prob}
Complete a prova de \olref[fol][com][mod]{lem:truth}.
\end{prob}

\tagendprob

\tagprob[notFOL]{probOr,probAnd,probIf}

\begin{prob}
Complete a prova de \olref[pl][com][mod]{lem:truth}.
\end{prob}

\tagendprob

\end{document}
